% Cahier des charges réorganisé - Queue Wait Time Estimation
\documentclass[11pt,a4paper]{article}
\usepackage[utf8]{inputenc}
\usepackage[T1]{fontenc}
\usepackage[french]{babel}
\usepackage{geometry}
\usepackage{hyperref}
\geometry{margin=25mm}
\title{Cahier des charges — Système d'estimation du temps d'attente aux caisses}
\author{Stage PFE — Salmen Ben Ammar}
\date{Version révisée -- 26 février 2026}

\begin{document}
\maketitle
\vspace{5mm}

\section{Contexte et définition du problème}
La surveillance des files d'attente en zone caisses à l'aide de caméras de surveillance permet d'optimiser les ressources opérationnelles (ouverture de caisses, appel d'agents) sans dépendre des données internes des caisses.

Le problème adressé est l'estimation fiable, en temps réel, de la longueur des files, du débit de service et du temps d'attente uniquement à partir d'observations vidéo, en tenant compte des incertitudes liées aux détections, aux occlusions et à la variabilité opérationnelle.

\section{Objectif de projet}
Fournir un système autonome capable de :
\begin{itemize}
  \item détecter et suivre les clients dans des zones d'intérêt configurables ;
  \item estimer le nombre de personnes en file, le débit de service et le temps d'attente en quasi‑temps réel ;
  \item quantifier l'incertitude des estimations ;
  \item générer des alertes et recommandations opérationnelles basées sur les estimations et leur fiabilité ;
  \item exposer les indicateurs via un tableau de bord et permettre l'historisation et l'export.
\end{itemize}

\section{Acteurs}
\begin{itemize}
  \item \textbf{Superviseur / Responsable magasin} : réceptionne alertes, décide ouverture de caisses.
  \item \textbf{Opérateur système} : configuration des zones, seuils et maintenance du système.
  \item \textbf{Système (automatisé)} : détecte, calcule métriques, émet alertes et stocke historiques.
\end{itemize}

\section{Les actions associées à chaque acteur}
\begin{itemize}
  \item Superviseur : consulter dashboard, valider ou rejeter recommandations, planifier ressources.
  \item Opérateur système : définir polygones de zones, configurer seuils/persistance, surveiller logs et qualité de détection.
  \item Système : générer événements (arrivée, sortie), calculer taux et temps d'attente, estimer incertitudes, déclencher webhooks/notifications.
\end{itemize}

% Force sections 5-7 to start on page 2
\clearpage
\section{Choix de technologie}
Technologies effectivement utilisées dans le projet (extrait de `requirements.txt` et du code) :
\begin{itemize}
  \item Langage : Python 3.x
  \item Détection : Ultralytics YOLO26 via le package `ultralytics` (modèles `yolo26*.pt`).
  \item Traitement d'image : `opencv-python` (OpenCV) et `supervision` pour aides géométriques/tracking.
  \item Tracking / helpers : `supervision` (SORT/associations) et logique interne dans `src/tracker.py`.
  \item Calculs et data : `numpy`, `pandas`.
  \item GUI : `tkinter` (interface locale) et `Pillow` (PIL) pour gestion d'images.
  \item Intégration & webhooks : `requests` pour envoi d'alertes HTTP.
  \item Outils dev / visualisation : `jupyter`, `matplotlib`.
  \item Stockage : exports CSV (via `csv_logger`), option SQLite pour historisation.
\end{itemize}

\paragraph{Site web / API}
Le projet comprendra un backend REST et un dashboard web. Backend recommandé : \textbf{FastAPI} + `uvicorn` exposant endpoints (ex. `/metrics`, `/alerts`, `/history`, `/zones`) et un webhook ; WebSocket/SSE pour mises à jour temps réel. Frontend : React/Vue (production) ou \textbf{Streamlit} (prototype).

Dépendances web requises : `fastapi`, `uvicorn`, `websockets`/`python-socketio` et (optionnel) toolchain frontend.

\section{Module à développer}
La solution se compose des modules suivants :
\begin{itemize}
  \item \texttt{detector} : intégration du modèle de détection et prétraitement des frames.
  \item \texttt{tracker} : association des détections et maintien des identifiants persistants.
  \item \texttt{zone\_manager} : définition et gestion des zones d'intérêt, détection des transitions.
  \item \texttt{queue\\_analyzer} : comptage, estimation des taux (arrivée/service) et calcul du temps d'attente.
  \item \texttt{uncertainty} : propagation des incertitudes et production d'intervalles ou scores de fiabilité.
  \item \texttt{webhook} / \texttt{notifier} : règles d'alerte et envoi de notifications/executions d'actions.
  \item \texttt{gui} / \texttt{dashboard} : affichage en temps réel et historisation.
  \item \texttt{api} / \texttt{backend} : API REST et services temps réel (WebSocket/SSE) exposant métriques, historiques et contrôles.
  \item \texttt{frontend} : application web (React/Vue) ou prototype Streamlit consommant l'API pour supervision distante.
  \item \texttt{csv\\_logger} : export et stockage des métriques pour analyse offline.
\end{itemize}

\section{Méthodologie de conception adoptée}
Approche incrémentale et modulaire : développement par sprints courts, tests sur vidéos d'évaluation et validation empirique.

Principes adoptés :
\begin{itemize}
  \item séparation claire des responsabilités (détection, tracking, analyse, interface) ;
  \item pipeline temps réel avec traitements par fenêtres temporelles pour lissage et stabilisation ;
  \item quantification d'incertitude depuis les entrées (incertitude de détection, perte d'observabilité) par propagation analytique ou Monte‑Carlo simplifié ;
  \item règles d'alerte paramétrables et testables automatiquement.
\end{itemize}

% Force architecture to start on page 3
\clearpage
\section{Architecture}
L'architecture proposée comprend :
\begin{itemize}
  \item Acquisition vidéo (RTSP/MP4) \rightarrow Prétraitement des frames.
  \item Module de détection \rightarrow Module de tracking \rightarrow Zone Manager (génération d'événements).
  \item Queue Analyzer et Uncertainty Engine qui calculent indicateurs et intervalles.
  \item Notifier/Rules Engine qui déclenche alertes (webhook/email) et stocke métriques via CSV/DB.
  \item Dashboard (GUI) qui consomme API pour affichage temps réel et historiques.
\end{itemize}

Un diagramme d'architecture peut être ajouté en annexe (schéma bloc représentant les flux entre modules).

\section{Choix technologiques (détaillés)}
\begin{itemize}
  \item Modèle de détection : utiliser un modèle quantique‑prêt pour inférence CPU/GPU selon disponibilité (YOLO family).
  \item Tracking : SORT pour simplicité/latence, DeepSORT si ré‑identification requise.
  \item Backend web/API : \textbf{FastAPI} + `uvicorn` pour endpoints REST et WebSocket/SSE ; `python-socketio` ou `websockets` si nécessaire.
  \item Base de données : SQLite pour preuve de concept; Postgres pour déploiement multi‑instance.
  \item Observabilité : logs structurés (JSON), métriques simples (latence, fps, taux d'erreur) exposées via endpoint.
  \item Déploiement : containerisation possible (Docker), exécution locale pour prototypes.
\end{itemize}

\vfill
\noindent PDF généré depuis le dépôt "Queue-Wait-Time-Estimation" — version réorganisée.

\end{document}
